\chapter*{前言：从“视觉组教程”到机器人算法教材}
\addcontentsline{toc}{chapter}{前言：从“视觉组教程”到机器人算法教材}

旧版教程诞生于装甲板检测占据视觉组主要研发资源的时期。它记录了当时完整而真实的学习路径，也因此积累了大量目标检测模型、损失函数、图像滤波和工程技巧。然而，机器人算法的目标不是在离线数据集上得到一个框或一张掩膜，而是在传感器、算力、通信、执行器和赛事规则共同约束下，持续估计环境并产生可执行决策。

新版以闭环系统为主线。设机器人真实状态为 $\vect{x}_k$，传感器产生量测 $\vect{y}_k$，算法估计状态 $\hat{\vect{x}}_k$，决策与规划生成参考 $\vect{r}_k$，控制器输出指令 $\vect{u}_k$：
\begin{align}
  \vect{x}_{k+1} &= f(\vect{x}_k,\vect{u}_k,\vect{w}_k), \\
  \vect{y}_k &= h(\vect{x}_k)+\vect{v}_k, \\
  \hat{\vect{x}}_k &= \mathcal{E}(\vect{y}_{0:k},\vect{u}_{0:k-1}), \\
  \vect{r}_k &= \mathcal{P}(\hat{\vect{x}}_k,\mathcal{M}_k), \\
  \vect{u}_k &= \mathcal{C}(\hat{\vect{x}}_k,\vect{r}_k).
\end{align}
这里 $\vect{w}_k$ 和 $\vect{v}_k$ 分别表示过程扰动与量测噪声，$\mathcal{M}_k$ 表示地图、任务和规则等上下文。全书每一部分都对应这条链上的一个映射，并最终在自瞄、哨兵与雷达站案例中重新闭合。

这也改变了内容取舍。传统图像处理只保留理解数字图像、卷积、特征点和多视图几何所必需的内容；目标检测从模型史缩减为任务形式、数据闭环、指标、现代基线与部署；Kalman 滤波等动态状态估计则被保留并加强，因为它们承担的是“从带噪量测恢复动态状态”的核心工作，而非简单平滑图像。

本稿仍处于结构重构阶段。稳定的数学和工程原理与易变化的赛季规则分开维护；涉及 2025/2026 赛季的事实将注明规则版本、发布日期和来源。旧稿不会立即删除，而是作为迁移证据保留，直到新版相应章节达到可替代的质量。

